%------------------------------------------------------------------------------%

\usepackage{calculo_varias_variaveis-1}

\title {Máximos e Mínimos}

%------------------------------------------------------------------------------%
\begin{document}

\addtitle

%------------------------------------------------------------------------------%
\section{Extremos Locais}

%------------------------------------------------------------------------------%
%\framedgraphic{Derivada Direcional}{derivada_direcional}{}

%------------------------------------------------------------------------------%
\begin{frame}{Extremos Locais}
  Seja $f(x, y)$ definida em uma região $R$ que contém o ponto $(a, b)$

  \pause
  \begin{itemize}[<+->]
    \item \(f(a, b)\) é um valor \emph{máximo local} de $f$ se
          \[
            f(a, b) \geq f(x, y)
          \]
          para \((x, y)\) do domínio em um disco aberto centrado em \((a, b)\)
    \item \(f(a, b)\) é um valor \emph{mínimo local} de $f$ se
          \[
            f(a, b) \leq f(x, y)
          \]
          para \((x, y)\) do domínio em um disco aberto centrado em \((a, b)\)
  \end{itemize}
\end{frame}

%------------------------------------------------------------------------------%
\begin{frame}{Teste da Derivada de Primeira Ordem}
  Se $f(x, y)$ tiver um máximo ou mínimo local em um ponto interior \((a, b)\)
  do seu domínio

  \pause
  \vspace{\baselineskip}
  e se as derivadas de primerira ordem existirem em \((a, b)\), então

  \pause
  \vspace{\baselineskip}
  \[
    f_x(a, b) = 0
    \qquad
    f_y(a, b) = 0
  \]
\end{frame}

%------------------------------------------------------------------------------%
\begin{frame}{Ponto Crítico}
  Um ponto interior do domínio de uma função \(f(x, y)\) onde

  \pause
  \vspace{\baselineskip}
  \(f_x\) e \(f_y\) sejam zero, ou

  \pause
  \vspace{\baselineskip}
  \(f_x\) ou \(f_y\) não existam

  \pause
  \vspace{\baselineskip}
  é um \emph{Ponto Crítico} de $f$
\end{frame}

%------------------------------------------------------------------------------%
\begin{frame}{Ponto de Sela}
  Uma função diferenciável \(f(x, y)\) possui um \emph{ponto de sela} em \((a, b)\)

  \pause
  \vspace{\baselineskip}
  se em todo disco aberto centrado em \((a, b)\)

  \pause
  \vspace{\baselineskip}
  existirem pontos \((x, y)\) onde
  \[
    f(x, y) > f(a, b)
  \]
  \pause
  e existirem pontos \((x, y)\) onde
  \[
    f(x, y) < f(a, b)
  \]
\end{frame}

%------------------------------------------------------------------------------%
\addtocounter{exemplo}{1}
\begin{frame}{Exemplo \theexemplo}
  Encontre os valores extremos locais de
  \[
     f(x, y) = x^2 + y^2 - 4y + 9
  \]
\end{frame}

\begin{frame}{Exemplo \theexemplo}
  Derivadas parciais
  \pause
  \vspace{\baselineskip}
  \[
    \D{f}{x}
    \pause
    = \D{}{x}\left(x^2 + y^2 - 4y + 9\right)
    \pause
    = 2x
  \]
  \pause
  \vspace{\baselineskip}
  \[
    \D{f}{y}
    \pause
    = \D{}{y}\left(x^2 + y^2 - 4y + 9\right)
    \pause
    = 2y -4
  \]
\end{frame}

\begin{frame}{Exemplo \theexemplo}
  A função está definida no plano todo

  \pause
  \vspace{\baselineskip}
  Pontos críticos

  \pause
  \vspace{\baselineskip}
  As derivadas existem em todos os pontos do plano

  \pause
  \vspace{\baselineskip}
  Plano tangente horizontal
  \[
    f_x(x, y) = 2x = 0
    \pause
    \quad\Rightarrow\quad
    \pause
    x = 0
  \]
  \[
    \pause
    f_y(x, y) = 2y -4 = 0
    \pause
    \quad\Rightarrow\quad
    \pause
    y = 2
  \]
\end{frame}

\begin{frame}{Exemplo \theexemplo}
  O único ponto critico é \((0, 2)\)

  \pause
  \vspace{\baselineskip}
  Onde a função vale
  \[
    f(0, 2)
    \pause
    = \left(x^2 + y^2 - 4y + 9\right)\evalat{(0, 2)}{}
    \pause
    = 0 + 4 - 8 + 9
    \pause
    = 5
  \]

  \pause
  \vspace{\baselineskip}
  Analizando a função percebemos que 5 é o menor valor que ela assume,
  portanto é um ponto de mínimo local (e nesse caso global também)
\end{frame}

%------------------------------------------------------------------------------%
\section{Teste da Derivada Segunda}

%------------------------------------------------------------------------------%
\begin{frame}{Hessiana}
  Matriz Hessiana
  \[
    \arraycolsep=2mm
    \def\arraystretch{1.2}
    H(x, y) = \begin{bmatrix}
      \null\;f_{xx}(x, y)\;\null & \null\;f_{xy}(x, y)\;\null \\
      f_{yx}(x, y) & f_{yy}(x, y)
    \end{bmatrix}
  \]

  \pause
  \vspace{\baselineskip}
  Discrimiante (ou hessiano)
  \[
    \arraycolsep=2mm
    \def\arraystretch{1.2}
    \det H(x, y) \;=\; f_{xx}(x, y)\,f_{yy}(x, y) \;-\; f_{xy}^2(x, y)
  \]
\end{frame}

%------------------------------------------------------------------------------%
\begin{frame}{Teste da Derivada Segunda}
  Seja $f(x, y)$ diferenciável e
  \(
    f_x(a, b) \;=\; f_y(a, b) \;=\; 0
  \)

  \pause
  \vspace{\baselineskip}
  \begin{enumerate}[<+->]
    \item $f$ tem um \emph{máximo local} em $(a, b)$ se
    \[
      f_{xx}(a, b) \bm{<} 0
      \quad\text{e}\quad
      f_{xx}(a, b) f_{yy}(a, b) - f^2_{xy}(a, b) \bm{>} 0
    \]
    \item $f$ tem um \emph{mínimo local} em $(a, b)$ se
    \[
      f_{xx}(a, b) \bm{>} 0
      \quad\text{e}\quad
      f_{xx}(a, b) f_{yy}(a, b) - f^2_{xy}(a, b) \bm{>} 0
    \]
  \end{enumerate}
\end{frame}

%------------------------------------------------------------------------------%
\begin{frame}{Teste da Derivada Segunda}
\begin{enumerate}[<+->]
  \setcounter{enumi}{2}
  \item $f$ tem um \emph{ponto de sela} em $(a, b)$ se
  \[
    f_{xx}(a, b) f_{yy}(a, b) - f^2_{xy}(a, b) \bm{<} 0
  \]
  \item o teste é inconclusivo se
  \[
    f_{xx}(a, b) f_{yy}(a, b) - f^2_{xy}(a, b) \bm{=} 0
  \]
\end{enumerate}
\end{frame}

%------------------------------------------------------------------------------%
\addtocounter{exemplo}{1}
\begin{frame}{Exemplo \theexemplo}
  Encontre os valores extremos locais da função
  \[
    f(x, y) = xy -x^2 - y^2 -2x -2y +4
  \]
\end{frame}

\begin{frame}{Exemplo \theexemplo}
  Derivadas parciais
  \pause
  \vspace{\baselineskip}
  \[
    \D{f}{x}
    \pause
    = \D{}{x}\left(xy -x^2 - y^2 -2x -2y +4\right)
    \pause
    = y -2x -2
  \]
  \pause
  \vspace{\baselineskip}
  \[
    \D{f}{y}
    \pause
    = \D{}{y}\left(xy -x^2 - y^2 -2x -2y +4\right)
    \pause
    = x -2y -2
  \]
\end{frame}

\begin{frame}{Exemplo \theexemplo}
  A função está definida no plano todo

  \pause
  \vspace{\baselineskip}
  Pontos críticos

  \pause
  \vspace{\baselineskip}
  As derivadas existem em todos os pontos do plano
\end{frame}

\begin{frame}{Exemplo \theexemplo}
\begin{columns}
\column{0.5\linewidth}
  \begin{align*}
    \onslide<+->{f_x(x, y) & = y -2x -2 = 0 \\}
    \onslide<+->{f_y(x, y) & = x -2y -2 = 0}
  \end{align*}

  \vspace{\baselineskip}
  \begin{align*}
    \onslide<+->{x -2y -2 & = 0      \\}
    \onslide<+->{x        & = 2y + 2}
  \end{align*}

\column{0.5\linewidth}
  \begin{align*}
    \onslide<+->{y -2x -2         & = 0  \\}
    \onslide<+->{y -2(2y + 2) - 2 & = 0  \\}
    \onslide<+->{y -4y + -6       & = 0  \\}
    \onslide<+->{-3y              & = 6  \\}
    \onslide<+->{y                & = -2}
  \end{align*}

  \vspace{\baselineskip}
  \[
    \onslide<+->{x
    = 2y + 2}
    \onslide<+->{= 2(-2) + 2}
    \onslide<+->{= -2}
  \]
\end{columns}
\end{frame}

\begin{frame}{Exemplo \theexemplo}
  Calculando as derivadas segundas
  \[
    \pause
    \DS{f}{x}
    \pause
    = \D{}{x}\left(y -2x -2\right)
    \pause
    = -2
  \]
  \[
    \pause
    \DS{f}{y}
    \pause
    = \D{}{y}\left(x -2y -2\right)
    \pause
    = -2
  \]
  \[
    \pause
    \DM{f}{x}{y}
    \pause
    = \D{}{x}\left(x -2y -2\right)
    \pause
    = 1
  \]

  \pause
  Discriminante
  \[
    \pause
    D(x, y)
    \pause
    = f_{xx}(-2, -2) f_{yy}(-2, -2) - f^2_{xy}(-2, -2)
    \pause
    = (-2)(-2) - 1
    \pause
    = 3
  \]
\end{frame}

\begin{frame}{Exemplo \theexemplo}
  Como \(f_{xx}(-2, -2) = -2 < 0\)
  \pause
  e \(D=3>0\)

  \pause
  \vspace{\baselineskip}
  O ponto \((-2, -2)\) é um ponto de mínimo local

  \pause
  \vspace{\baselineskip}
  O valor mínimo da função é
  \[
    f(-2, -2)
    \pause
    = \left(xy -x^2 - y^2 -2x -2y +4\right)\evalat{(-2, -2)}{}
  \]
  \[
    \hphantom{f(-2, -2)}
    \pause
    = (-2)(-2) - (-2)^2 - (-2)^2 -2(-2) -2(-2) + 4
  \]
  \[
    \hphantom{f(-2, -2)}
    \pause
    = 4 -4 -4 +4 +4 + 4
  \]
  \[
    \hphantom{f(-2, -2)}
    \pause
    = 8
  \]
\end{frame}

%------------------------------------------------------------------------------%
\addtocounter{exemplo}{1}
\begin{frame}{Exemplo \theexemplo}
  Encontre os valores extremos locais da função
  \[
    f(x, y) = 3y^2 -2y^3 -3x^2 +6xy
  \]
\end{frame}

\begin{frame}{Exemplo \theexemplo}
  Derivadas parciais
  \pause
  \vspace{\baselineskip}
  \[
    \D{f}{x}
    \pause
    = \D{}{x}\left(3y^2 -2y^3 -3x^2 +6xy\right)
    \pause
    = -6x +6y
    \pause
    = 6y -6x
  \]
  \pause
  \vspace{\baselineskip}
  \[
    \D{f}{y}
    \pause
    = \D{}{y}\left(3y^2 -2y^3 -3x^2 +6xy\right)
    \pause
    = 6y -6y^2 +6x
  \]
\end{frame}

\begin{frame}{Exemplo \theexemplo}
  A função está definida no plano todo

  \pause
  \vspace{\baselineskip}
  Pontos críticos

  \pause
  \vspace{\baselineskip}
  As derivadas existem em todos os pontos do plano
\end{frame}

\begin{frame}{Exemplo \theexemplo}
\begin{columns}
\column{0.5\linewidth}
  \begin{align*}
    \onslide<+->{f_x(x, y) & = 6y -6x = 0 \\}
    \onslide<+->{f_y(x, y) & = 6y -6y^2 +6x = 0}
  \end{align*}
  \begin{align*}
    \onslide<+->{6y -6x & = 0 \\}
    \onslide<+->{x      & = y}
  \end{align*}
  \column{0.5\linewidth}
  \begin{align*}
    \onslide<+->{6y -6y^2 +6x & = 0 \\}
    \onslide<+->{6y^2 -6y -6x & = 0 \\}
    \onslide<+->{6y^2 -6y -6y & = 0 \\}
    \onslide<+->{6y^2 - 12y   & = 0 \\}
    \onslide<+->{6y(y-2)      & = 0}
  \end{align*}

  \vspace{\baselineskip}
  \onslide<+->{Portanto \(y=0\) ou \(y=2\)}

  \vspace{\baselineskip}
  \onslide<+->{Os pontos críticos são \((0, 0)\) e \((2, 2)\)}
\end{columns}
\end{frame}

\begin{frame}{Exemplo \theexemplo}
  Calculando as derivadas segundas
  \[
    \pause
    \DS{f}{x}
    \pause
    = \D{}{x}\left(6y -6x\right)
    \pause
    = -6
  \]
  \[
    \pause
    \DS{f}{y}
    \pause
    = \D{}{y}\left(6y -6y^2 +6x\right)
    \pause
    = 6 -12y
  \]
  \[
    \pause
    \DM{f}{x}{y}
    \pause
    = \D{}{y}\left(6y -6x\right)
    \pause
    = 6
  \]
\end{frame}

\begin{frame}{Exemplo \theexemplo}
  \onslide<+->{Discriminante}
  \begin{align*}
    \onslide<+->{D(x, y)}
    \onslide<+->{& = f_{xx}(0, 0) f_{yy}(0, 0) - f^2_{xy}(0, 0) \\}
    \onslide<+->{& = (-6)(6-12y) - 6^2 \\}
    \onslide<+->{& = -36 +72y - 36 \\}
    \onslide<+->{& = 72(y-1)}
  \end{align*}
\end{frame}

\begin{frame}{Exemplo \theexemplo}
  Analizando o ponto \((0, 0)\)

  \pause
  \vspace{\baselineskip}
  Como
  \[
    \pause
    D(0, 0)
    \pause
    = 72(0-1)
    \pause
    = -72
    \pause
    \;<\; 0
  \]
  \pause
  esse ponto é um ponto de sela
\end{frame}

\begin{frame}{Exemplo \theexemplo}
  \onslide<+->{Analizando o ponto \((2, 2)\)}

  \onslide<+->{Como}
  \[
    \onslide<+->{D(0, 0) = 72(2-1) = 72 > 0}
    \qquad
    \onslide<+->{f_{xx} = -6 < 0}
  \]
  \onslide<+->{o ponto \((2, 2)\) é um máximo local com valor}
  \begin{align*}
    \onslide<+->{f(2, 2)
    & = \left(3y^2 -2y^3 -3x^2 +6xy\right)\evalat{(2, 2)}{} \\}
    \onslide<+->{& = 3\times 2^2 -2 \times 2^3 -3\times 2^2 + 6\times 2 \times 2\\}
    \onslide<+->{& = 12 -16 -12 + 24 \\}
    \onslide<+->{& = 8}
  \end{align*}
\end{frame}

%------------------------------------------------------------------------------%
\section{Lista Mínima}

%------------------------------------------------------------------------------%
\begin{frame}{Lista Mínima}
  Cálculo Vol. 2 do Thomas 12\textsuperscript{a} ed. -- Seção 14.7
  \begin{enumerate}
    \item Estudar o texto da seção
    \item Resolver os exercícios: 2, 9, 11, 21, 23, 24, 25, 27
  \end{enumerate}

  \vfill
  Atenção: A prova é baseada no livro, não nas apresentações
\end{frame}

%------------------------------------------------------------------------------%
\end{document}
%------------------------------------------------------------------------------%
